% MAS1004 Assignment 2: short report, English template.
%
% Write freely about what you learned from this assignment:
%   - what you tried, and why
%   - how the assignment went
%   - what you learned
%
% Write the text on one page, yourself, without an LLM. After the text you
% may add up to three pages of figures.
% Compile with XeLaTeX (Overleaf: Menu, Compiler, XeLaTeX).

\documentclass[11pt]{article}
\usepackage[onepage]{mas1004}

\reporttitle{Assignment 2: Short report}
\studentname{TRUONG HOANG HAI ANH}
\studentid{20231165}

\begin{document}
\makeheader

I tried:
I ran the starter code for both tracks without changing anything. For Track A, I trained the starter recipe on all 44,257 training photos, two times with different seeds. This was to see how much the score changes just by luck. For Track B, I trained the starter ResNet18 for 4 epochs and checked that the model passes the size limits. I did not remove or add any photos.

The results were:
Track A seed 0 = 0.4965.
Track A seed 1 = 0.4924.
Track B = 0.4987.

I was surprised that:
Just changing the seed moved the score by about 0.004. That is bigger than the gap between Track B and Track A seed 0 (0.0023). I was also surprised that the cities are so different. Tokyo is easy (about 0.60 to 0.64), but Osaka is hard (about 0.34 to 0.37). Many Osaka photos were guessed as Tokyo.

I learned that:
A small score difference does not always mean one run is better, because luck can change the score too. I also learned that one overall score hides a lot. The score for each city and the confusion table show much more about what the model gets wrong.

If I had more time:
I would look at the photos the model got wrong with high confidence, especially Osaka photos guessed as Tokyo. Then I would try one change at a time, for example removing indoor photos, and compare with the starter result.

Overall:
I learned that before trying to make a model better, I should first check how much its score changes by luck.

\figurepages

\end{document}
